\title{DeepSeekMath: Pushing the Limits of Mathematical Reasoning in Open Language Models}

\begin{abstract}
Language models often struggle with mathematical reasoning due to its inherent complexity and structured requirements. This work presents DeepSeekMath~7B, developed by further pre-training DeepSeek-Coder-Base-v1.5 7B with 120B math-centric tokens, extracted from Common Crawl along with supplemental natural language and code resources. DeepSeekMath~7B attains a remarkable 51.7\% on the MATH benchmark at the competition level without leveraging ensemble methods or external tools, coming close to the capabilities demonstrated by models such as Gemini-Ultra and GPT-4. When evaluated using self-consistency with 64 samples, DeepSeekMath~7B secures 60.9\% on MATH. The advanced mathematical reasoning demonstrated by DeepSeekMath~is grounded in two principal components: Firstly, our carefully curated data selection pipeline unlocks the latent potential in publicly available web datasets. Secondly, we propose Group Relative Policy Optimization (GRPO), an enhancement of Proximal Policy Optimization (PPO) that both strengthens mathematical reasoning and improves PPO's memory efficiency.
\end{abstract}

\section{Introduction}

Large language models (LLMs) have fundamentally transformed methods of mathematical reasoning within artificial intelligence, leading to notable progress on quantitative reasoning assessments \citep{MATH} as well as geometry reasoning evaluations \citep{trinh2024solving}. In addition to their impact on automated problem-solving, these models have also demonstrated their value in supporting humans tackling intricate mathematical challenges \citep{tao}. Despite these achievements, state-of-the-art systems like GPT-4 \citep{gpt4} and Gemini-Ultra \citep{gemini} remain unavailable to the public, while existing open-source alternatives exhibit markedly inferior performance.

This paper presents DeepSeekMath, a specialized language model developed to enhance mathematical problem-solving capabilities, significantly surpassing those of existing open-source systems and approaching GPT-4’s performance across standard academic evaluations. Central to our approach is the construction of the DeepSeekMath Corpus, a high-caliber pre-training dataset composed of 120B math-specific tokens. This resource is assembled from the Common Crawl (CC) via a classifier built on the fastText framework \citep{joulin2016fasttext}. For the initial phase, the classifier receives positive training data from OpenWebMath \citep{openwebmath}, while negative samples are drawn from a broad array of unrelated web pages. Utilizing this setup, the classifier extracts additional positive samples from the CC, which undergo subsequent curation and validation by human annotators. The classifier is iteratively retrained with this improved corpus, further elevating its selection accuracy. Empirical assessment demonstrates that our expansive corpus yields high-quality training material: DeepSeekMath-Base 7B attains scores of 64.2\% on GSM8K \citep{gsm8k} and 36.2\% on the advanced MATH benchmark \citep{MATH}, exceeding the performance of Minerva 540B \citep{minerva}. Moreover, as the DeepSeekMath Corpus encompasses multiple languages, we observe considerable performance gains on Chinese-language math benchmarks \citep{wei2023cmath,agieval}. We propose that our methodology for mathematical data collection sets a promising precedent for future research endeavors and anticipate continued advancements in this area.

DeepSeekMath-Base is built upon DeepSeek-Coder-Base-v1.5 7B \citep{deepseek-coder}, as we find that leveraging a model pre-trained on code tasks offers advantages over using a generic LLM for initialization. Additionally, our results reveal that mathematical instruction tuning brings performance gains not only for math-specific tasks but also improves outcomes on the MMLU \citep{mmlu} and BBH \citep{bbh} benchmarks, suggesting that such training bolsters both mathematical proficiency and broader reasoning skills in the model.

Following the pre-training phase, we conduct mathematical instruction tuning on DeepSeekMath-Base utilizing datasets focused on chain-of-thought \citep{cot}, program-of-thought \citep{pot,pal}, and tool-integrated reasoning \citep{tora}. The produced model, DeepSeekMath-Instruct 7B, surpasses other models with 7B parameters and achieves performance levels similar to instruction-tuned models of 70B parameters available in the open-source community.

In addition, we propose Group Relative Policy Optimization (GRPO), a modified reinforcement learning (RL) approach inspired by Proximal Policy Optimization (PPO) \citep{schulman2017proximal}. Unlike PPO, GRPO omits the use of a critic model and instead derives its baseline from aggregated group scores, thereby substantially lowering the computational costs required for training. When trained only on a select subset of English instruction tuning examples, GRPO achieves marked performance gains over the robust DeepSeekMath-Instruct model, enhancing results across both familiar benchmarks (GSM8K: 82.9\% $\rightarrow$ 88.2\%, MATH: 46.8\% $\rightarrow$ 51.7\%) as well as in generalization to out-of-domain mathematical problems (for instance, CMATH: 84.6\% $\rightarrow$ 88.8\%) during the RL training phase. Furthermore, we articulate a unified framework that connects methodologies such as Rejection Sampling Fine-Tuning (RFT) \citep{yuan2023scaling}, Direct Preference Optimization (DPO) \citep{dpo}, PPO, and GRPO. Through this lens, all these strategies can be interpreted as variations or simplifications of core RL concepts. To systematically analyze the critical properties of this framework, we conduct comprehensive evaluations encompassing factors like online versus offline learning, the effects of supervising outcomes compared to processes, and the distinctions between single-step and iterative RL settings. Finally, we elucidate how our RL approach enhances instruction-tuned model performance and discuss promising avenues for refining RL effectiveness within this cohesive paradigm.

\subsection{Contributions}
We present scalable approaches to math pre-training, complemented by an investigation and evaluation of reinforcement learning methods.

\textbf{Math Pre-Training at Scale}

\begin{itemize}[topsep=0pt]
    \item Through our investigation, we demonstrate that the Common Crawl dataset, which is publicly available, holds significant potential for mathematical research. Utilizing a carefully crafted selection methodology, we developed the DeepSeekMath Corpus—a comprehensive dataset encompassing 120B tokens from web sources identified as rich in mathematical content. This corpus notably surpasses the scale of math-oriented web pages employed in Minerva \citep{minerva} by nearly a factor of 7, and is also 9 times larger than the newly launched OpenWebMath \citep{openwebmath}.

\item DeepSeekMath-Base 7B, our base model trained from scratch, matches the performance of Minerva 540B \citep{minerva}, suggesting that model size alone does not dictate proficiency in mathematical reasoning. Notably, compact models exposed to high-quality training data can also exhibit robust performance.

\item We present our results from experiments involving mathematical training. Pretraining models with code prior to mathematics increases their proficiency in solving math problems, whether or not tools are used. These findings contribute to the ongoing debate surrounding the question: \textit{does code training enhance reasoning skills?} Our results suggest that it does—specifically for mathematical reasoning tasks.

\item Despite the widespread use of arXiv papers as training data, particularly within mathematics-focused research, this approach yields no significant gains across any of the mathematical benchmarks evaluated in this study.
\end{itemize}

\textbf{Exploration and Analysis of Reinforcement Learning}

\begin{itemize}[topsep=0pt]
    \item We present Group Relative Policy Optimization (GRPO), a novel reinforcement learning method that achieves both efficiency and effectiveness. Unlike Proximal Policy Optimization (PPO), GRPO eliminates the need for a critic model by utilizing group scores to estimate the baseline, thereby lowering the computational demands required for training.
    \item Our results show that GRPO substantially improves the performance of the instruction-tuned DeepSeekMath-Instruct model using only instruction-tuning data. Additionally, we find that this approach leads to better out-of-domain generalization during the reinforcement learning phase.
    \item We introduce a comprehensive framework that encapsulates various reinforcement learning techniques including RFT, DPO, PPO, and GRPO. We also perform thorough empirical studies—examining aspects like online versus offline training, outcome versus process supervision, and single-turn versus iterative reinforcement learning—to systematically analyze the crucial components of this framework.
    \item Building upon this unified framework, we investigate the mechanisms underlying the efficacy of reinforcement learning, and identify several promising avenues for advancing the reinforcement learning capabilities of large language models (LLMs).
\end{itemize}

\subsection{Summary of Evaluations and Metrics}

\begin{itemize}[topsep=0pt]
    \item \textbf{English and Chinese Mathematical Reasoning}: Our evaluation systematically examines model performance on both English and Chinese mathematical reasoning tasks, with benchmarks ranging from elementary to collegiate levels. For English, the assessments utilize GSM8K \citep{gsm8k}, MATH \citep{MATH}, SAT \citep{llemma}, OCW Courses \citep{minerva}, and MMLU-STEM \citep{mmlu}. The Chinese evaluation suite comprises MGSM-zh \citep{mgsm}, CMATH \citep{wei2023cmath}, Gaokao-MathCloze \citep{agieval}, and Gaokao-MathQA \citep{agieval}. We investigate the models' capabilities in producing complete, text-based solutions independently of external tools, as well as their proficiency in solving mathematical problems through Python.

On English evaluation benchmarks, DeepSeekMath-Base demonstrates performance on par with the proprietary Minerva 540B \citep{minerva} and consistently outperforms all available open-source base models, including Mistral 7B \citep{mistral} and Llemma-34B \citep{llemma}, regardless of their exposure to mathematical pre-training, frequently by a wide margin. Remarkably, DeepSeekMath-Base achieves even greater results on Chinese benchmarks. This can likely be attributed to the model’s data collection strategy, which diverges from previous approaches \citep{minerva, llemma} by incorporating high-quality math datasets in multiple languages rather than restricting to English-only data. Furthermore, with subsequent mathematical instruction fine-tuning and reinforcement learning, both DeepSeekMath-Instruct and DeepSeekMath-RL attain strong outcomes, surpassing 50\% accuracy on the challenging, competition-level MATH dataset—a first for any open-source model.

\item \textbf{Formal Mathematics}: DeepSeekMath-Base is assessed on its ability to translate informal mathematical statements into formal proofs through the task outlined by \citep{dsp_proof}, employing Isabelle \citep{isabelle} as the proof assistant within the miniF2F \citep{minif2f} dataset. The results reveal that DeepSeekMath-Base achieves impressive few-shot autoformalization outcomes.
\item \textbf{Natural Language Understanding, Reasoning, and Code}: In order to gain a detailed understanding of the model’s general proficiency in comprehension, logical reasoning, and programming, we conduct evaluations of DeepSeekMath-Base on the Massive Multitask Language Understanding (MMLU) benchmark \citep{mmlu}, which features 57 multiple-choice examinations spanning a wide range of disciplines. Additionally, we utilize BIG-Bench Hard (BBH) \citep{bbh}—comprising 23 advanced tasks that mainly demand complex, multi-step reasoning—and the HumanEval \citep{codex} and MBPP \citep{mbpp} datasets that are standard for code generation assessment. Results indicate that mathematical pre-training enhances performance in both language understanding and reasoning capabilities.

\section{Math Pre-Training}

\subsection{Data Collection and Decontamination} This section details the methodology employed to build the DeepSeekMath Corpus utilizing data sourced from Common Crawl. As illustrated in Figure \ref{fig:data_collect}, we introduce a stepwise pipeline that facilitates the comprehensive extraction of a substantial mathematical text corpus from Common Crawl, beginning with an initial seed dataset (such as a concise, high-quality set of mathematics-related texts). It should be emphasized that this framework is adaptable and can similarly be applied in other areas, including programming.

Initially, we select OpenWebMath \citep{openwebmath}—a curated dataset of high-quality mathematical web documents—as our foundational seed corpus. With this dataset, we proceed to train a fastText model \citep{joulin2016fasttext} to help identify additional web pages with characteristics similar to those found in OpenWebMath. For this process, we randomly extract 500,000 samples from the seed corpus to serve as positive examples, while another 500,000 web pages are sampled from Common Crawl to function as negative examples. The training utilizes a publicly available fastText implementation\footnote{\url{https://fasttext.cc}}, with hyperparameters set as follows: vector size of 256, learning rate at 0.1, a maximum word n-gram length of 3, a minimum threshold of 3 for word occurrences, and 3 training epochs. To reduce the massive scale of Common Crawl, we implement deduplication at the URL level as well as near-duplicate removal, trimming the dataset to 40 billion HTML documents. We deploy the fastText model to retrieve mathematical pages from this deduplicated web dataset. In order to further screen for high-caliber mathematical content, we rank these retrieved web pages by their fastText-assigned scores and retain only the highest-ranked items. The quantity of preserved material is determined based on outcomes from pre-training experiments on data volumes of 40B, 80B, 120B, and 160B tokens. For this initial cycle, we opt to retain the top 40B tokens.

Following the initial round of data collection, a significant number of mathematical web pages remain missing. This shortfall primarily arises because the fastText model was originally trained on a positive example set with limited variety. To enhance the performance of the model, we expand the seed corpus by incorporating more mathematical web sources. Our method involves segmenting the entire Common Crawl dataset into non-overlapping domains, where each domain is defined by web pages under a common base URL. We then compute, for each domain, the proportion of its web pages that were collected in the first iteration. Domains with more than 10\% of their pages gathered in this round are considered to be math-oriented (such as \url{mathoverflow.net}).

For these math-oriented domains, we undertake manual annotation to identify URLs that are specifically linked to mathematical content (for instance, \url{mathoverflow.net/questions}). Any mathematical web pages under these URLs that have not yet been collected are incorporated into our seed set. By increasing the number and range of positive examples in this way, the fastText model is retrained, which significantly enhances its ability to retrieve mathematical data in the next iteration. After conducting four rounds of data collection, our process yields a total of 35.5M mathematical web pages, comprising 120B tokens. At the fourth stage, we observe that almost 98\% of these pages had already been retrieved in the third iteration, prompting us to terminate further data collection.

In order to prevent contamination of our benchmarks, we adopt the approach outlined in \cite{deepseek-coder}, excluding any web content that features questions or answers from established English mathematical benchmarks such as GSM8K \citep{gsm8k} and MATH \citep{MATH}, as well as Chinese resources like CMATH \citep{wei2023cmath} and AGIEval \citep{agieval}. Specifically, our filtering mechanism removes any segment of text in the math training dataset that contains a 10-gram sequence matching any substring within the referenced evaluation benchmarks. Additionally, for benchmark materials with lengths shorter than 10 grams but consisting of at least 3 grams, we utilize an exact string matching process to detect and eliminate potentially contaminated web pages.

\subsection{Validating the Quality of the DeepSeekMath~Corpus}

We run pre-training experiments to investigate how the DeepSeekMath~Corpus is compared with the recently released math-training corpora:

\begin{itemize}[topsep=0pt]
    \item \textbf{MathPile} \citep{mathpile}: a composite dataset encompassing 8.9B tokens derived from multiple sources such as textbooks, Wikipedia, ProofWiki, CommonCrawl, StackExchange, and arXiv, with arXiv accounting for more than 85\% of the material;
    \item \textbf{OpenWebMath} \citep{openwebmath}: data extracted from CommonCrawl and screened specifically for mathematical relevance, amounting to 13.6B tokens;
    \item \textbf{Proof-Pile-2} \citep{llemma}: a specialized mathematical dataset formed by combining OpenWebMath, AlgebraicStack (containing 10.3B tokens of mathematical code), and 28.0B tokens from arXiv papers. For experiments involving Proof-Pile-2, we adhere to the arXiv:Web:Code token distribution of 2:4:1, as prescribed by \cite{llemma}.
\end{itemize}

\subsubsection{Training Setting}
\label{sec:quality-policy}
We fine-tune a foundational language model with 1.3B parameters, which utilizes the same architecture as the DeepSeek LLMs \citep{deepseek-llm}, referenced as DeepSeek-LLM 1.3B. The model is individually fine-tuned on each mathematical dataset for a total of 150B tokens. All training is executed using the efficient and streamlined HAI-LLM \citep{haillm} framework. In line with the methodologies adopted for DeepSeek LLMs, we employ the AdamW optimizer \citep{adamW} with parameters $\beta_1=0.9$, $\beta_2=0.95$, and a $\mathrm{weight\_decay}$ of 0.1. The learning rate schedule follows a multi-step approach: the learning rate is gradually increased to its peak value after 2,000 warmup iterations, reduced to 31.6\% of the maximum by the 80\% mark of total training steps, and finally lowered to 10.0\% of its maximum value at the 90\% point. The learning rate’s upper bound is set at 5.3e-4, while training proceeds with a batch size of 4M tokens and a context length capped at 4K tokens.

\subsubsection{Evaluation Results}

\textbf{The DeepSeekMath~Corpus is of high quality, covers multilingual mathematical content, and is the largest in size.}

\begin{itemize}[topsep=0pt]
    \item \textbf{High-quality}:
    We assess downstream capabilities across eight mathematical tasks using few-shot chain-of-thought prompting \cite{cot}. Table \ref{tab:corpora_comparison} demonstrates a distinct advantage in performance for models trained on the DeepSeekMath~Corpus. Furthermore, Figure \ref{fig:corpora_comparisons} illustrates that, after exposure to 50B tokens (equivalent to a complete epoch of Proof-Pile-2), models trained with DeepSeekMath Corpus consistently outperform those using Proof-Pile-2, underscoring the superior average quality of DeepSeekMath Corpus data.
    \item \textbf{Multilingual}:
    The DeepSeekMath~Corpus contains multilingual content, with a strong emphasis on English and Chinese, which are the predominant languages represented. Table \ref{tab:corpora_comparison} confirms that training with DeepSeekMath~Corpus yields notable gains in mathematical reasoning abilities for both English and Chinese. By contrast, other mathematical corpora—mostly centered on English—offer limited benefit or may even impair performance when dealing with Chinese mathematical reasoning tasks.
    \item \textbf{Large-scale}:
    In terms of scale, the DeepSeekMath~Corpus is significantly larger compared to previously available mathematical corpora. Figure \ref{fig:corpora_comparisons} reveals that DeepSeek-LLM 1.3B, when trained on DeepSeekMath~Corpus, achieves more rapid and persistent improvement, reflected in a steeper learning trajectory. On the other hand, baseline corpora are considerably smaller and subjected to repeated epochs during training, which leads to model performance stagnating after a short period.
\end{itemize}

\subsection{Training and Evaluating DeepSeekMath-Base 7B}

In this section, we present DeepSeekMath-Base 7B, a foundational model that demonstrates advanced capabilities in mathematical reasoning. The model’s initialization is based on DeepSeek-Coder-Base-v1.5 7B \citep{deepseek-coder}, followed by training on 500B tokens. The composition of the training corpus is distributed as follows: 56\% DeepSeekMath Corpus, 4\% AlgebraicStack, 10\% arXiv sources, 20\% Github code, with the remaining 10\% comprising natural language text from Common Crawl in both English and Chinese. The training procedure primarily adheres to the configuration in Section \ref{sec:quality-policy}; however, we adjust the maximum learning rate to 4.2e-4 and utilize a batch size of 10M tokens.

We present an extensive evaluation of DeepSeekMath-Base 7B, particularly examining its proficiency in generating fully self-contained mathematical solutions independently, its aptitude for solving mathematical tasks utilizing external tools, and its effectiveness in formal theorem proving. In addition to mathematics, we offer a broader characterization of the base model by assessing its capabilities in natural language comprehension, logical reasoning, and programming.

\paragraph{Mathematical Problem Solving with Step-by-Step Reasoning}

We assess the capabilities of DeepSeekMath-Base in addressing mathematical problems via few-shot chain-of-thought prompting \citep{cot}, utilizing eight distinct benchmarks in both English and Chinese languages. These benchmarks span various aspects of quantitative reasoning—such as GSM8K \citep{gsm8k}, MATH \citep{MATH}, and CMATH \citep{wei2023cmath}—as well as multiple-choice question formats—including MMLU-STEM \citep{mmlu} and Gaokao-MathQA \citep{agieval}. The evaluation covers a broad spectrum of mathematical disciplines, ranging from elementary through collegiate levels of difficulty.

As detailed in Table \ref{tab:base_math_cot}, among all open-source base models—including the prominent Mistral 7B \citep{mistral} and the newer Llemma 34B \citep{llemma}, which has been trained on Proof-Pile-2 \citep{llemma} for mathematics—DeepSeekMath-Base 7B consistently achieves the highest scores across all eight benchmark tasks. Particularly on the challenging, competition-caliber MATH dataset, DeepSeekMath-Base attains more than a 10\% absolute gain over other open-source base models. Furthermore, it exceeds the performance of Minerva 540B \citep{minerva}—a proprietary model based on PaLM \citep{palm}, enhanced with math data, and possessing 77 times more parameters.

\paragraph{Mathematical Problem Solving with Tool Use} Our assessment of program-assisted mathematical reasoning leverages the GSM8K and MATH benchmarks with few-shot program-of-thought prompting \citep{pot,pal}. For each task, the models generate Python code solutions, incorporating libraries like \textit{math} and \textit{sympy} as needed for complex calculations. The final answer is determined by executing the generated code. According to Table \ref{tab:base_math_pot_proof}, DeepSeekMath-Base 7B demonstrates superior performance compared to the previously leading Llemma 34B model.

\paragraph{Formal Mathematics}

The automation of formal proofs plays a crucial role in guaranteeing the precision and dependability of mathematical reasoning while also improving productivity, a trend that has seen growing interest lately. In this context, we assess DeepSeekMath-Base 7B using the informal-to-formal proving task introduced in \citep{dsp_proof}, where the objective is to construct a formal proof from an informal statement, its formal equivalent, and an accompanying informal proof. Our evaluation utilizes the miniF2F benchmark \citep{minif2f}, which focuses on Olympiad-level mathematics problems that require formalization. For each problem, we use few-shot prompting to generate an Isabelle formal proof. Consistent with the methodology in \cite{dsp_proof}, we employ language models to produce proof sketches and subsequently apply the standard automated prover Sledgehammer \citep{sledgehammer} to complete the remaining proof steps. Table \ref{tab:base_math_pot_proof} highlights that DeepSeekMath-Base 7B achieves notable results in the task of proof autoformalization.

\paragraph{Natural Language Understanding, Reasoning, and Code}

We assess the ability of the models in natural language understanding using MMLU \citep{mmlu}, for reasoning with BBH \citep{bbh}, and for programming proficiency through HumanEval \citep{codex} and MBPP \citep{mbpp}. According to the results displayed in Table \ref{tab:understanding_reasoning_code}, DeepSeekMath-Base 7B achieves notable gains in both MMLU and BBH scores when compared to its predecessor, DeepSeek-Coder-Base-v1.5 \citep{deepseek-coder}, highlighting the advantageous effects of mathematical training on capabilities in language understanding and reasoning. Furthermore, by incorporating code tokens during continued training, DeepSeekMath-Base 7B successfully preserves the coding performance achieved by DeepSeek-Coder-Base-v1.5 on the same two coding tasks. In summary, DeepSeekMath-Base 7B demonstrates a substantial performance advantage over the general-purpose Mistral 7B \citep{mistral} across the evaluated reasoning and coding tasks.

\section{Supervised Fine-Tuning}

\subsection{SFT Data Curation}

We develop a comprehensive instruction-tuning dataset for mathematics, encompassing both English and Chinese questions from a range of mathematical disciplines and difficulty levels. Each problem is matched with its corresponding solution, provided in chain-of-thought (CoT) \citep{cot}, program-of-thought (PoT) \citep{pot,pal}, and tool-integrated reasoning styles \citep{tora}. The dataset includes a total of 776K training samples.

\begin{itemize}[topsep=0pt]
    \item \textbf{English mathematical datasets}: Our dataset includes GSM8K and MATH, each supplemented with solutions that incorporate tool-based reasoning. Additionally, we utilize a curated portion of MathInstruct \citep{MathInstruct}, and integrate the training examples from Lila-OOD \citep{lila}, where the solutions employ either CoT or PoT methodologies. The assembled English datasets encompass a broad spectrum of mathematical domains, including but not limited to algebra, probability, number theory, calculus, and geometry.
    \item \textbf{Chinese mathematical datasets}: We compile a range of Chinese K-12 mathematics problems, covering 76 distinct sub-categories such as linear equations. For each problem, we provide solution annotations utilizing both CoT and tool-integrated reasoning approaches.
\end{itemize}

\subsection{Training and Evaluating DeepSeekMath-Instruct 7B}

In this section, we present DeepSeekMath-Instruct 7B, which is developed by applying mathematical instruction tuning to DeepSeekMath-Base. To construct training inputs, samples are randomly combined until the maximum context window of 4K tokens is filled. The model is trained over 500 steps using a batch size of 256, maintaining a fixed learning rate of 5e-5 throughout the process.

We evaluate models' mathematical performance both without and with tool use, on 4 quantitative reasoning benchmarks in English and Chinese. We benchmark our model against the leading models of the time:

\begin{itemize}[topsep=0pt]
    \item \textbf{Closed-source models} encompass the following: (1) the GPT lineup, with GPT-4 \citep{gpt4} and GPT-4 Code Interpreter~\footnote{\url{https://openai.com/blog/chatgpt-plugins\#\#code-interpreter}} representing the most advanced variants, (2) Gemini Ultra and Pro \citep{gemini}, (3) Inflection-2 \citep{inflection-2}, (4) Grok-1~\footnote{\url{https://x.ai/model-card}}, and, in addition, newly introduced models by Chinese enterprises, including (5) Baichuan-3~\footnote{\url{https://www.baichuan-ai.com}}, and (6) the newest GLM-4~\footnote{\url{https://open.bigmodel.cn/dev/api\#glm-4}}

    from the GLM family \citep{glm}.

These models are designed for general applications, with the majority having been subjected to various alignment processes.  
\item \textbf{Open-source models} comprise: general-purpose models such as (1) DeepSeek-LLM-Chat 67B \citep{deepseek-llm}, (2) Qwen 72B \citep{qwen}, (3) SeaLLM-v2 7B \citep{seallm}, and (4) ChatGLM3 6B \citep{chatglm3}. Additionally, several models offer advancements in mathematical reasoning, including (5) InternLM2-Math 20B~\footnote{\url{https://github.com/InternLM/InternLM-Math}}, which extends InternLM2 by first undergoing mathematics-specific training and subsequently instruction tuning, (6) Math-Shepherd-Mistral 7B, which employs PPO-based optimization \citep{schulman2017proximal} with a process-supervised reward structure applied to Mistral 7B \citep{mistral}, (7) WizardMath series \citep{wizardmath}, enhancing the mathematical capabilities of both Mistral 7B and Llama-2 70B \citep{llama2} using evolve-instruct (instruction tuning utilizing AI-generated instructions) and PPO techniques with tasks primarily from GSM8K and MATH datasets, (8) MetaMath 70B \citep{metamath}, representing Llama-2 70B fine-tuned on expanded versions of GSM8K and MATH, (9) ToRA 34B \cite{tora}, where CodeLlama 34B is further refined for mathematical problem solving that incorporates tools, and (10) MAmmoTH 70B \citep{MathInstruct}, a variant of Llama-2 70B subjected to instruction tuning utilizing MathInstruct.

As indicated in Table \ref{tab:sft_rl_math}, when tool use is restricted, DeepSeekMath-Instruct 7B exhibits robust capabilities in sequential reasoning tasks. Remarkably, on the competition-level MATH dataset, our model outperforms all other open-source systems and most proprietary alternatives (including Inflection-2 and Gemini Pro) by a margin of at least 9\% absolute. This advantage persists even against much larger models (such as Qwen 72B) and those augmented through reinforcement learning with a math specialization (like WizardMath-v1.1 7B). Although DeepSeekMath-Instruct achieves results comparable to the Chinese closed-source models GLM-4 and Baichuan-3 on MATH, it remains behind GPT-4 and Gemini Ultra.

When assessed in a framework that permits both natural language reasoning and the utilization of programmatic tools to tackle problems, DeepSeekMath-Instruct 7B attains nearly 60\% accuracy on the MATH dataset, outperforming all previously available open-source models. Additionally, across other evaluation benchmarks, our model demonstrates performance comparable to DeepSeek-LLM-Chat 67B, the former leading model in this space, despite being one-tenth its size.

\subsection{Group Relative Policy Optimization} Following the Supervised Fine-Tuning (SFT) phase, reinforcement learning (RL) has demonstrated its capability in enhancing large language models' (LLMs) skills in mathematical reasoning~\citep{wang2023math,wizardmath}. Here, we present Group Relative Policy Optimization (GRPO), our novel RL approach that is designed to be both efficient and robust.

\subsubsection{From PPO to GRPO} Proximal Policy Optimization (PPO) \citep{schulman2017proximal} is an actor-critic RL algorithm that is widely used in the RL fine-tuning stage of LLMs \citep{ouyang2022training}. In particular, it optimizes LLMs by maximizing the following surrogate objective:

\begin{equation}
\footnotesize
    \mathcal{J}_{PPO}(\theta) = \mathbb{E}{[q \sim P(Q),\ o \sim \pi_{\theta_{old}}(O|q)]} \frac{1}{|o|} \sum_{t=1}^{|o|} \min \left[ \frac{\pi_\theta(o_{t} | q, o_{<t})}{\pi_{\theta_{old}}(o_{t} | q, o_{<t})} A_{t}, \text{clip} \left( \frac{\pi_\theta(o_{t} | q, o_{<t})}{\pi_{\theta_{old}}(o_{t} | q, o_{<t})}, 1 - \varepsilon, 1 + \varepsilon \right)  A_{t} \right] ,
\end{equation}

The PPO objective, $\mathcal{J}_{PPO}(\theta)$, computes the expectation over queries $q$ sampled from $P(Q)$ and outputs $o$ generated from the old policy $\pi_{\theta_{old}}(O|q)$. For each token position $t$ in the sequence $o$, the objective averages over the minimum between the policy ratio multiplied by the advantage term, $A_{t}$, and the advantage weighted by the clipped ratio, where the ratio is bounded between $1 - \varepsilon$ and $1 + \varepsilon$. This formulation ensures that the policy update remains within a trust region by preventing large deviations from the previous policy.

Here, $\pi_{\theta}$ denotes the updated policy model while $\pi_{\theta_{old}}$ refers to its previous iteration. The variables $q$ and $o$ stand for questions drawn from a dataset and outputs produced by the earlier policy $\pi_{\theta_{old}}$. The parameter $\varepsilon$ acts as a hyper-parameter for clipping, as proposed in PPO, ensuring greater stability during optimization. The term $A_t$ represents the advantage, which is calculated using Generalized Advantage Estimation (GAE) \citep{gae}, leveraging the sequence of rewards $\{r_{\ge t}\}$ along with the predicted values from a trained value function $V_{\psi}$. Accordingly, PPO requires the concurrent training of a value function and the policy itself. To counter excessive reward model optimization, a commonly adopted technique is to include a token-wise KL-divergence penalty from a fixed reference model in the reward calculation at every generation step \citep{ouyang2022training}, formally described as

\begin{equation}
    r_{t} = r_\varphi(q, o_{\le t}) - \beta \log\frac{\pi_{\theta}(o_{t}|q, o_{<t})}{\pi_{ref}(o_{t}|q, o_{<t})},
\label{eq:PPO-reward}
\end{equation}

Here, $r_\varphi$ denotes the reward model, while $\pi_{ref}$ represents the reference model—typically, this is the original SFT model. The parameter $\beta$ serves as the scaling factor for the KL divergence penalty.

As the value function employed in PPO is typically another model of comparable size as the policy model, it brings a substantial memory and computational burden. Additionally, during RL training, the value function is treated as a baseline in the calculation of the advantage for variance reduction. While in the LLM context, usually only the last token is assigned a reward score by the reward model, which may complicate the training of a value function that is accurate at each token. To address this, as shown in Figure \ref{fig:grpo}, we propose Group Relative Policy Optimization (GRPO), which obviates the need for additional value function approximation as in PPO, and instead uses the average reward of multiple sampled outputs, produced in response to the same question, as the baseline. More specifically, for each question $q$, GRPO samples a group of outputs $\{o_1, o_2, \cdots, o_G\}$  from the old policy  $\pi_{\theta_{old}}$  and then optimizes the policy model by maximizing the following objective:

\begin{equation}
\footnotesize
\begin{split}
    \mathcal{J}_{GRPO}(\theta) &= \mathbb{E}_{q \sim P(Q), \{o_i\}_{i=1}^G \sim \pi_{\theta_{old}}(O|q)}  \\
    & \frac{1}{G}\sum_{i=1}^G \frac{1}{|o_i|} \sum_{t=1}^{|o_i|} \left\{ \min \left[ \frac{\pi_\theta(o_{i,t} | q, o_{i,<t})}{\pi_{\theta_{old}}(o_{i,t} | q, o_{i,<t})} \hat{A}_{i,t},\ \text{clip} \left( \frac{\pi_\theta(o_{i,t} | q, o_{i,<t})}{\pi_{\theta_{old}}(o_{i,t} | q, o_{i,<t})}, 1 - \varepsilon, 1 + \varepsilon \right)  \hat{A}_{i,t} \right] - \beta \mathbb{D}_{KL}\left[\pi_{\theta} || \pi_{ref}\right]\right\} ,
\end{split}
\label{eq:GRPO-obj}
\end{equation}

where $\varepsilon$ and $\beta$ are hyper-parameters, and $\hat{A}_{i,t}$  is the advantage calculated based on relative rewards of the outputs inside each group only, which will be detailed in the following subsections. The group relative way that GRPO leverages to calculate the advantages, aligns well with the comparative nature of rewards models,  as reward models are typically trained on datasets of comparisons between outputs on the same question. Also note that, instead of adding KL penalty in the reward, GRPO regularizes by directly adding the KL divergence between the trained policy and the reference policy to the loss, avoiding complicating the calculation of  $\hat{A}_{i,t}$. And different from the KL penalty term used in (\ref{eq:PPO-reward}), we estimate the KL divergence with the following unbiased estimator \citep{kl_approx}:

\begin{equation}
\small
    \mathbb{D}_{KL}\left[\pi_{\theta} || \pi_{ref}\right] = \frac{\pi_{ref}(o_{i,t}|q,o_{i,<t})}{\pi_{\theta}(o_{i,t}|q,o_{i,<t})}- \log\frac{\pi_{ref}(o_{i,t}|q,o_{i,<t})}{\pi_{\theta}(o_{i,t}|q,o_{i,<t})} - 1,
\end{equation}

which is guaranteed to be positive.

\begin{algorithm}[t]
  \small
  \caption{Iterative Group Relative Policy Optimization}
  \textbf{Input} initial policy model $\pi_{\theta_{\text{init}}}$; reward models $r_\varphi$; task prompts $\mathcal{D}$; 
  hyperparameters $\varepsilon$, $\beta$, $\mu$
  \begin{algorithmic}[1]
    \State Initialize policy model: $\pi_\theta \gets \pi_{\theta_{\text{init}}}$
    \For{iteration = 1, \dots, I}
       \State Set the reference model: $\pi_{ref} \gets \pi_{\theta}$
      \For{step = 1, \dots, M}
      \State Draw a mini-batch $\mathcal{D}_b$ from $\mathcal{D}$
      \State Store the current policy: $\pi_{\theta_{old}} \gets \pi_{\theta}$ 
      \State For each question $q$ in $\mathcal{D}_b$, generate $G$ samples $\{o_i\}_{i=1}^G$ from $\pi_{\theta_{old}} (\cdot \mid q)$
      \State Calculate rewards $\{r_i\}_{i=1}^{G}$ for the generated outputs $o_i$ by evaluating with $r_\varphi$
      \State For each sampled output, determine $\hat{A}_{i,t}$ at the $t$-th token using group relative advantage estimation.
      \For{GRPO iteration = 1, \dots, $\mu$}
        \State Perform policy update for $\pi_{\theta}$ by maximizing the GRPO objective (Equation \ref{eq:GC-GRPO})
      \EndFor
    \EndFor 
    \State Continuously refine $r_\varphi$ employing a replay-based training approach.
    \EndFor 
  \end{algorithmic}
  \textbf{Output} $\pi_\theta$
  \label{alg:iter-grpo}
\end{algorithm}

\subsubsection{Outcome Supervision RL with GRPO} Specifically, for every question $q$, a set of candidate outputs $\{o_1, o_2, \cdots, o_G\}$ is generated using the previous policy $\pi_{\theta_{old}}$. Each output is evaluated by a reward model, resulting in a set of rewards $\mathbf{r}=\{r_1, r_2, \cdots, r_G\}$, one for each candidate. These raw rewards are then standardized within the group by subtracting the mean and dividing by the standard deviation of the set. Outcome supervision utilizes these standardized rewards, assigning each to the respective output $o_i$ as the final token-level reward. This normalized value serves as the advantage $\hat{A}_{i, t}$ for all tokens in $o_i$, given by $\hat{A}_{i, t} = \widetilde{r}_i = \frac{r_i- {\rm mean}(\mathbf{r})}{{\rm std}(\mathbf{r})}$. The policy model is subsequently updated to maximize the objective specified in equation (\ref{eq:GRPO-obj}).

\subsubsection{Process Supervision RL with GRPO} In contrast to outcome supervision, which only provides feedback after the entire output is generated, process supervision offers more granular guidance by delivering a reward after each individual reasoning step—this approach is particularly advantageous for intricate mathematical problem-solving. Building upon the framework introduced by \cite{wang2023math}, we also incorporate process supervision, where a reward model assigns scores to each intermediate step. Specifically, for a given question $q$ and a collection of $G$ sampled outputs $\{o_1, o_2, \cdots, o_G\}$, the process reward model evaluates each reasoning step, producing a set of rewards structured as follows: $\mathbf{R} = \{ \{r_1^{index(1)},\cdots,r_1^{index(K_1)}\}, \cdots,  \{r_G^{index(1)},\cdots,r_G^{index(K_G)}\} \}$, where $index(j)$ indicates the position of the $j$-th step's terminal token, and $K_i$ represents the total number of steps in the $i$-th output sequence. The raw rewards are subsequently standardized using the dataset-wide mean and standard deviation: $\widetilde{r}_i^{index(j)} = \frac{r_i^{index(j)} - {\rm mean(\mathbf{R})}}{{\rm std(\mathbf{R})}}$.
With these normalized stepwise rewards, the process supervision framework defines the advantage for each token as the cumulative sum of normalized rewards from all subsequent steps, i.e., $\hat{A}_{i, t} = \sum_{index(j) \ge t} \widetilde{r}_i^{index(j)}$.
The policy is then updated by maximizing the objective articulated in equation (\ref{eq:GRPO-obj}).

\subsubsection{Iterative RL with GRPO} Over time, as reinforcement learning advances, the original reward model may no longer provide adequate guidance for the updated policy model. To address this limitation, we investigate an iterative approach to RL using GRPO. In this paradigm, detailed in Algorithm \ref{alg:iter-grpo}, new datasets for the reward model are created from samples generated by the current policy model. The reward model is further improved through continual training, utilizing a replay strategy that integrates 10\% of previous samples. Following this, the existing policy model is designated as the reference, and further training of the policy model proceeds using the updated reward model.

\subsection{Training and Evaluating DeepSeekMath-RL}

Our reinforcement learning experiments utilize DeepSeekMath-Instruct 7B as the policy model. The RL training dataset comprises approximately 144K chain-of-thought questions drawn from the GSM8K and MATH subsets of the SFT data, specifically omitting other SFT questions to isolate the effects of RL on benchmarks for which training examples are not present during RL. The reward model training set is prepared as described in \citep{wang2023math}. We initialize the reward model using DeepSeekMath-Base 7B and adopt a learning rate of $2 \times 10^{-5}$. For GRPO optimization, the policy model’s learning rate is set to $1 \times 10^{-6}$, with a KL regularization coefficient fixed at $0.04$. We generate $64$ sampled outputs for each question, use a maximum sequence length of $1024$, and employ a batch size of $1024$ for training. The policy network receives a single update following each exploration episode. DeepSeekMath-RL 7B is evaluated on the same benchmarks as DeepSeekMath-Instruct 7B. For evaluation, GSM8K and MATH under chain-of-thought settings are treated as in-domain benchmarks for DeepSeekMath-RL 7B, whereas all other benchmarks are considered out-of-domain.

Table \ref{tab:sft_rl_math} presents a comparison of open- and closed-source models employing both chain-of-thought and tool-integrated reasoning on benchmarks in English and Chinese. Our observations are as follows: 1) DeepSeekMath-RL 7B achieves 88.2\% accuracy on GSM8K and 51.7\% accuracy on MATH when leveraging chain-of-thought reasoning. These results exceed those of any other open-source models within the 7B to 70B parameter range and outperform most closed-source models as well. 2) Notably, DeepSeekMath-RL 7B’s training is limited to chain-of-thought-format instruction tuning data from GSM8K and MATH, with the model initialized from DeepSeekMath-Instruct 7B. In spite of this restricted training data, DeepSeekMath-RL 7B consistently outperforms DeepSeekMath-Instruct 7B across every evaluated metric, underscoring the advantages provided by reinforcement learning.

\section{Discussion}
This section presents an overview of the insights gained from our experiments involving pre-training and reinforcement learning.

\subsection{Lessons Learnt in Pre-Training}

We begin by outlining our pre-training process. Except where noted, the training configurations employed correspond to those detailed in Section \ref{sec:quality-policy}. It is important to highlight that, in this section, the DeepSeekMath Corpus specifically refers to a dataset comprising 89 billion tokens, collected during the second round of data gathering.

\subsubsection{Code Training Benefits Mathematical Reasoning}

An often-cited but unproven theory posits that engaging in code training enhances reasoning abilities. In this work, we address this claim in part by focusing on mathematics: we demonstrate that code training leads to better mathematical reasoning performance in models, irrespective of the use of external tools.

To study how code training affects mathematical reasoning, we experimented with the following two-stage training and one-stage training settings:

\textbf{Two-Stage Training}

\begin{itemize}[topsep=0pt]
    \item \textbf{Code Training for 400B Tokens $\rightarrow$ Math Training for 150B Tokens}: We initially pre-train DeepSeek-LLM 1.3B on 400B tokens sourced from code, and subsequently fine-tune it using 150B tokens containing mathematical content;
    \item \textbf{General Training for 400B Tokens $\rightarrow$ Math Training for 150B Tokens}: For comparison, a baseline model is constructed by pre-training with 400B general-domain tokens—drawn from a large general text corpus developed by DeepSeek-AI—instead of code tokens. This setup is designed to elucidate whether code token pre-training provides superior benefits for downstream mathematical reasoning compared to general-domain token pre-training.
\end{itemize}

\textbf{One-Stage Training}

\begin{itemize}[topsep=0pt]
    \item \textbf{Math Pretraining with 150B Tokens}: DeepSeek-LLM 1.3B is pretrained exclusively on 150 billion tokens of mathematical data;
    \item \textbf{Joint Training with 400B Code Tokens and 150B Math Tokens}: Incorporating math fine-tuning after code training has been observed to reduce code-related capabilities. Therefore, we explore a joint training approach, blending code and math tokens in a single-phase training process, to determine if this not only enhances mathematical problem-solving but also reduces the risk of catastrophic forgetting.
\end{itemize}

\paragraph{Results}
Tables \ref{tab:code-math} and \ref{tab:code-math-general-eval} present the outcomes of downstream evaluations across varying training configurations.

Engaging in code training contributes positively to program-aided mathematical reasoning across both the two-stage and one-stage training paradigms. As indicated in Table \ref{tab:code-math}, within the two-stage training setup, exclusive code training already brings about a marked improvement in solving GSM8K and MATH tasks with Python. Subsequent math-focused training in the latter stage leads to additional gains. Notably, in the one-stage training scheme, interleaving code tokens with math tokens effectively counteracts the catastrophic forgetting observed in the two-stage approach, while also promoting advancements in both coding abilities (Table \ref{tab:code-math-general-eval}) and program-aided mathematical reasoning (Table \ref{tab:code-math}).

Training on code not only benefits mathematical reasoning in the absence of tool use but also demonstrates notable advantages in a two-stage training paradigm. In particular, the preliminary phase of code pretraining contributes to moderate gains and enhances the effectiveness of later math-specific training, which collectively culminate in superior overall results. By contrast, when code tokens and math tokens are merged in a single-stage curriculum, the model’s mathematical reasoning without tools is adversely affected. A plausible explanation is that DeepSeek-LLM 1.3B, given its relatively small model size, does not possess sufficient capacity to comprehensively learn both coding and mathematical skills simultaneously.

\subsubsection{ArXiv Papers Seem Ineffective in Improving Mathematical Reasoning}

ArXiv papers are commonly included as a component of math pre-training data \citep{minerva,gpt-f,llemma,mathpile}. However, detailed analysis regarding their impact on mathematical reasoning has not been extensively conducted. Perhaps counter-intuitively, according to our experiments, arXiv papers seem ineffective in improving mathematical reasoning. We experiment with models of different sizes, including DeepSeek-LLM 1.3B and DeepSeek-Coder-Base-v1.5 7B \citep{deepseek-coder}, using arXiv corpora that underwent varied processing pipelines:

\begin{itemize}[topsep=0pt]
    \item \textbf{MathPile} \citep{mathpile}: comprises 8.9 billion tokens, curated via a set of data cleaning and selection heuristics, with scientific papers from arXiv constituting more than 85\% of the dataset;
    \item \textbf{ArXiv-RedPajama} \citep{redpajama}: includes all arXiv LaTeX source files, where document preambles, comments, user-defined macros, and references have been stripped, resulting in a collection with a token count of 28.0 billion.
\end{itemize}

In our study, we independently train DeepSeek-LLM 1.3B for 150B tokens and DeepSeek-Coder-Base-v1.5 7B for 40B tokens using individual arXiv corpora. Our findings indicate that exposure to arXiv papers does not significantly enhance mathematical reasoning capabilities. When either model is trained solely on arXiv data, performance does not improve and sometimes even declines across a range of mathematical evaluation benchmarks with varying levels of difficulty featured in this work. These benchmarks encompass quantitative reasoning tasks such as GSM8K and MATH (Table \ref{tab:arxiv-cot}), multiple-choice assessments such as MMLU-STEM (Table \ref{tab:arxiv-cot}), and formal mathematics benchmarks, including miniF2F (Table \ref{tab:arxiv_atp}).

However, this conclusion has its limitations and should be taken with a grain of salt. We have not yet studied:

\begin{itemize}[topsep=0pt]
    \item The influence of arXiv tokens on additional mathematics-focused applications not explored in this study, like the process of translating formal mathematical theorems or proofs into more informal descriptions;
    \item The consequences of integrating arXiv tokens with diverse data sources;
    \item The extent to which the advantages provided by arXiv papers would persist or become apparent in models with greater capacity.
\end{itemize}

Thus, further exploration is required, which we leave for future studies.

\subsection{Insights of Reinforcement Learning}
\subsubsection{Towards to a Unified Paradigm} Here, we introduce a unified framework to systematically examine various training strategies, including SFT, RFT, DPO, PPO, and GRPO. We also perform empirical studies to investigate the key elements influencing this framework. In most cases, the gradient of a training algorithm with respect to the parameter $\theta$ is formulated as follows:

\begin{equation}
    \nabla_{\theta}\mathcal{J}_{\textcolor{red}{\mathcal{A}}}(\theta) = \mathbb{E}[\underbrace{(q,o) \sim \textcolor{red}{\mathcal{D}}}_{Data \ Source}]\left( \frac{1}{|o|} \sum_{t=1}^{|o|}  \underbrace{GC_{{\mathcal{A}}}(q, o, t, \textcolor{red}{\pi_{{rf}}})}_{Gradient \ Coefficient}  \nabla_{\theta}\log \pi_{\theta}(o_t | q, o_{<t})\right).
\label{eq:objective}
\end{equation}

There exist three key components: 1) \textit{Data Source $\mathcal{D}$}, which determines the training data; 2) \textit{Reward Function $\pi_{{rf}}$}, which is the source of the training reward signal; 3) \textit{Algorithm $\mathcal{A}$}: which processes the training data and the reward signal to the gradient coefficient $GC$ that determines the magnitude of the penalty or reinforcement for the data. We analyze several representative methods based on such a unified paradigm:

\begin{itemize}
    \item \textbf{Supervised Fine-tuning (SFT)}: This method adjusts a pretrained model by further training it on a dataset specifically curated by humans for SFT.
    \item \textbf{Rejection Sampling Fine-tuning (RFT)}: RFT enhances the SFT model through additional training on outputs generated by the SFT model in response to SFT prompts, selectively retaining only those outputs deemed correct.
    \item \textbf{Direct Preference Optimization (DPO)}: DPO builds upon the SFT model by applying fine-tuning using enriched outputs from the SFT model, employing pair-wise DPO loss to optimize preferences.
    \item \textbf{Online Rejection Sampling Fine-tuning (Online RFT)}: In contrast to RFT, Online RFT uses the SFT model to initialize the policy model, and then continually improves it by fine-tuning with outputs generated in real time by the current policy model.
    \item \textbf{PPO/GRPO}: In this approach, the SFT model serves as the starting point for the policy model, which is then strengthened by training on live outputs sampled from the ongoing policy model.
\end{itemize}

Table \ref{tab:data_policy} provides an overview of the components comprising these methods. For a comprehensive explanation of the derivation procedure, consult Appendix \ref{app:analysis-rl}.

\paragraph{Observation about Data Source} We categorize the data source into online sampling and offline sampling. Online sampling refers to obtaining training data directly from the ongoing explorations of the current policy model during training. In contrast, offline sampling involves collecting training data based on the outputs generated by the initial SFT model. Methods such as RFT and DPO utilize the offline approach, whereas Online RFT and GRPO employ the online approach.

As illustrated in Figure \ref{fig:rl-analysis}, our results indicate that Online RFT yields much stronger performance than RFT across both benchmarks. In the early phases of training, the performance of Online RFT and RFT remains similar, but as training advances, Online RFT obtains a pronounced edge, highlighting the benefits of the online training methodology. This observation aligns with expectations: at the outset, the actor’s behavior closely mirrors that of the SFT model, leading to only subtle distinctions in the data collected. As training proceeds, however, the samples generated by the actor begin to diverge more notably, thereby amplifying the advantages conferred by real-time data collection.

\paragraph{Observation about Gradient Coefficient} Within the algorithm, the reward signal is transformed into a gradient coefficient used to modify the model parameters. In our experimental setup, we categorize the reward function into `Rule' and `Model' components. The `Rule' is designed to evaluate response quality strictly according to answer correctness, whereas the `Model' involves training a reward model that assigns a score to each response; this reward model is itself trained using judgments based on the Rule criterion. As shown in Equations \ref{eq:GC-RFT} and \ref{eq:GC-GRPO}, GRPO differs fundamentally from Online RFT in its treatment of the gradient coefficient: GRPO uniquely modulates this coefficient according to the specific reward value obtained from the reward model. This feature enables GRPO to administer reinforcement or punishment in proportion to the response quality. On the other hand, Online RFT does not possess this capability; it does not penalize incorrect answers and delivers equal reinforcement to all correct responses, regardless of the quality or magnitude of correctness.

As illustrated in Figure \ref{fig:rl-analysis}, GRPO outperforms online RFT, which underscores the effectiveness of modifying positive and negative gradient coefficients. Moreover, GRPO+PS achieves better results than GRPO+OS, suggesting that incorporating step-aware, fine-grained gradient coefficients is advantageous. Additionally, we investigate the impact of iterative RL; our experiments involve two iterations. As depicted in Figure \ref{fig:iter-rl}, iterative RL leads to notable performance gains, with the most pronounced improvement occurring during the initial iteration.

\subsubsection{Why RL Works?} In this work, we apply reinforcement learning using a selected portion of instruction tuning data, which leads to substantial improvements over models trained solely with instruction tuning. To shed light on the underlying reasons for RL’s efficacy, we assess both Pass@K and Maj@K metrics for the Instruct and RL models on two different benchmarks. As depicted in Figure \ref{fig:maj-pass}, reinforcement learning boosts the Maj@K metric but does not substantially affect Pass@K. This suggests that RL’s advantage mainly stems from producing a more robust output distribution, specifically by elevating the likelihood of selecting the correct answer within the TopK responses, rather than by advancing the model’s core competence. In a similar vein, \citep{wang2023making} have documented a \textbf{misalignment problem} in SFT models for reasoning tasks, demonstrating that leveraging preference alignment strategies \citep{yuan2023rrhf,song2023preference,wang2023making} can lead to notable improvements in SFT model reasoning performance.

\subsubsection{How to Achieve More Effective RL?}
\label{sec:effec_rl}
Our findings indicate that RL performs effectively for tasks involving mathematical reasoning. Moreover, we introduce a cohesive framework that encompasses various representative training approaches, interpreting them through the lens of either explicit or streamlined RL methodologies. Within this framework, as detailed in Equation \ref{eq:objective}, the core elements comprise the Data Source, Algorithm, and Reward Function. We also outline several promising avenues for advancing research with respect to these three essential components.

\paragraph{Data Source} The data source constitutes the foundational input for all training methodologies. Within the scope of RL, we define the data source as unlabeled questions along with the outputs produced by sampling from the policy model. For this study, we restrict ourselves to employing questions provided during the instruction tuning phase and utilize straightforward nucleus sampling for generating responses. We hypothesize that this choice partially accounts for the RL pipeline's enhancements being limited to Maj@K metrics. In subsequent work, we intend to apply our RL pipeline to prompts that are outside the training distribution, while simultaneously leveraging \textbf{advanced sampling (decoding) strategies} such as tree-search-based techniques \citep{yao2023tree}. Furthermore, \textbf{efficient inference techniques} \citep{xia-etal-2023-speculative,leviathan2023fast,kwon2023efficient,xia2024unlocking}, which impact how effectively policy models can explore, are also recognized as highly significant in this context.

\paragraph{Algorithms} In these frameworks, algorithms utilize the available data and reward signals to compute the gradient coefficient used for model parameter updates. As indicated by Equation \ref{eq:objective}, prevailing approaches heavily \textbf{TRUST} the reward function’s output to determine adjustments in the conditional probability of a particular token. Nevertheless, it remains unattainable to guarantee the absolute reliability of the reward signal, particularly when tackling highly complex tasks. For instance, even in the PRM800K datasets \citep{lightman2023let}, where annotations are conducted meticulously by trained professionals, roughly 20\% of entries are still labeled incorrectly\footnote{\url{https://github.com/openai/prm800k/issues/12\#issuecomment-1728491852}}. With this in mind, we consider reinforcement learning algorithms capable of withstanding noisy or imperfect reward feedback. We contend that \textbf{WEAK-TO-STRONG} alignment strategies \citep{burns2023weak} represent a transformative shift in the landscape of learning algorithms.

\paragraph{Reward Function} The reward function serves as the main driver for learning in RL. In most cases, the reward function is realized through a neural reward model. We identify three key research areas concerning reward models: 1) \textbf{Improving the generalization capacity of reward models.} Reward models need to generalize robustly in order to deal with questions outside their training distribution and more advanced decoding scenarios; if they do not, RL may only serve to preserve the status quo of LLMs’ outputs instead of actually enhancing their abilities; 2) \textbf{Capturing and representing the reward model’s uncertainty.} Accurately quantifying uncertainty has the potential to mediate between less capable reward models and learning strategies that are intended to improve from weak to strong performance; 3) \textbf{Constructing high-quality process reward models in a data- and compute-efficient manner}, which can generate nuanced supervision signals during the reasoning process \citep{lightman2023let,wang2023math}.

\section{Conclusion, Limitation, and Future Work}

In this work, we introduce DeepSeekMath, which establishes new state-of-the-art results among open-source models on the MATH competition benchmark and achieves performance levels close to leading proprietary systems. DeepSeekMath is based on the DeepSeek-Coder-v1.5 7B architecture and is continuously trained over 500B tokens, with 120B of these tokens drawn from mathematical content identified within Common Crawl data. Through comprehensive ablation experiments, we identify that online resources, particularly web pages, serve as a valuable source of mathematical content, whereas arXiv appears to contribute less than initially anticipated. To further advance mathematical reasoning, we propose Group Relative Policy Optimization (GRPO), a modified approach based on Proximal Policy Optimization (PPO), capable of enhancing reasoning abilities while reducing memory requirements. Empirical findings demonstrate the efficacy of GRPO, even when DeepSeekMath-Instruct 7B already attains strong results on established benchmarks. Additionally, we present a unified framework for interpreting a variety of reinforcement learning methodologies and discuss several promising avenues for further improvement of RL techniques.

While DeepSeekMath~demonstrates strong results on quantitative reasoning tasks, its performance in areas such as geometry and theorem-proving lags behind that of proprietary models. For example, our preliminary experiments reveal the model's inability to address questions involving triangles and ellipses, which may be attributed to potential data selection biases present during the pre-training and fine-tuning stages. Moreover, owing to limitations in model size, DeepSeekMath~does not match GPT-4's effectiveness in few-shot learning scenarios. Whereas GPT-4 benefits significantly from few-shot prompts, DeepSeekMath~exhibits comparable outcomes under both zero-shot and few-shot settings. Moving forward, we intend to enhance our data selection pipeline to assemble a more robust pre-training dataset, as well as to investigate alternative strategies for more efficient reinforcement learning with large language models, as outlined in Section~\ref{sec:effec_rl}.

\end{document}